% Generated draft: requires mathematical and source review.
\documentclass[11pt]{article}
\usepackage[margin=25mm]{geometry}
\usepackage{fontspec}
\setmainfont{Latin Modern Roman}
\usepackage{amsmath,amssymb,mathtools,unicode-math}
\setmathfont{Latin Modern Math}
\usepackage{enumitem,needspace,xurl,hyperref}
\setlength{\parindent}{0pt}
\setlength{\parskip}{5pt}
\setlength{\emergencystretch}{3em}
\newcommand{\N}{\mathbb{N}}
\newcommand{\Z}{\mathbb{Z}}
\newcommand{\Q}{\mathbb{Q}}
\newcommand{\R}{\mathbb{R}}
\newcommand{\C}{\mathbb{C}}
\newcommand{\F}{\mathbb{F}}
\newcommand{\A}{\mathbb{A}}
\newcommand{\PP}{\mathbb{P}}
\newcommand{\E}{\mathbb{E}}
\newcommand{\eps}{\varepsilon}
\newcommand{\dd}{\,\mathrm{d}}
\begin{document}
\section*{UnsolvedMath ID: 30001337}
\textbf{Source problem code:} catalogue entry
\subsection{Definitions and mathematical statement}
A ring \(R\) has zero divisors if there exist \(a,b\in R\setminus\{0\}\) with \(ab=0\). The supplied record does not specify whether the rings under consideration are commutative or unital, or impose any additional hypotheses on them.

An element \(v\) of a \(C^*\)-algebra \(A\) is a partial isometry if \(vv^*v=v\). Its initial and final projections are, respectively,
\[
 p_v=v^*v,\qquad q_v=vv^*.
\]
These definitions explain possible meanings of the support projections mentioned in the source; the record does not specify whether initial projections, final projections, or both are intended.

\textbf{Construction problem.} For rings with zero divisors, find a mathematically reasonable organization of the support projections of partial isometries implementing ring multiplication in a ring \(C^*\)-algebra construction, so that the construction retains all relevant multiplicative information.

The source does not specify the ring \(C^*\)-algebra construction to be extended, the meaning of implementing multiplication, the required relations among the partial isometries and their projections, or a criterion for retaining the relevant information. Accordingly, this is a construction question with unresolved specifications, rather than a fully quantified existence assertion with a defined success criterion.
\subsection{Short English statement}
The problem asks how to organize the support projections arising when partial isometries represent multiplication in rings with zero divisors, while preserving the relevant multiplicative information. The intended ring class, operator relations, and information-preservation criterion remain unspecified.
\subsection{Sources}
\begin{itemize}
\item Local UnsolvedMath registry: supplied record with permanent ID \(30001337\), source problem code \texttt{catalogue entry}, title \emph{Ring \(C^*\)-Algebras with Zero Divisors}, in the Oberwolfach Reports Open Problems collection. The proposed year is recorded as 2009; no individual proposer is supplied.
\item Supplied source citation: \emph{Noncommutative Geometry (2010)}, Oberwolfach Reports, pp.\ 2271--2334, DOI \texttt{10.4171/owr/2009/41}. Original Oberwolfach report: \url{https://doi.org/10.4171/owr/2009/41}.
\item Supplied literature reference: \emph{K-theory of ring C-star-algebras}, \url{https://arxiv.org/abs/0906.4903}.
\item Supplied literature reference: \emph{K-theory erratum}, \url{https://arxiv.org/abs/1209.3486}.
\item Literature assessment dated \texttt{2026-08-21}: the supplied curation labels the extension to rings with zero divisors open and states that the integer case and its \(K\)-theory are solved. This is supplied dated curation, not an independently verified assessment.
\item The registry labels the statement \texttt{corrected\_verified} and reports that its mathematical content was checked against the official report, with notation, context, and question form normalized. That verification claim is supplied provenance and has not been independently checked here.
\end{itemize}
\end{document}
